% NAF 9544, batch 2 — Gallica views 21–40.
% Pass: Opus 5.5 (claude-opus-5-5), 2026-09-22 — first pass, unchecked against the views by a human.
%
% The volume. An album of autographs mounted on guards, each piece announced
% by a guard leaf bearing a large ink lot number, a pasted cutting from a
% sale catalogue and a note of the sale. The catalogue dates it 1601-1700,
% copied verbatim into \dating{} although this batch alone runs from 1620
% (Mersenne) through the Richelieu epigrams (c. 1642) to a Germain fragment
% (before 1831) and two Italian telegrams of 1884.
%
% What the twenty views hold.
%   v. 21–23  Sophie Germain: an untitled, undated, unsigned philosophical
%             fragment on « l'idée fondamentale d'une science », headed by a
%             later hand « [Ce qui suit s'est trouvé dans les papiers laissés
%             par Mlle Germain] ». Begins on v. 21 (a small leaf paginated 16),
%             continues on the left of v. 22 (its verso), then the right of
%             v. 22 (a second leaf paginated 17), and ends two lines into the
%             verso of that leaf (v. 23). Nothing follows it. The hand was
%             compared with her autograph at Français 9115 v. 222 and is
%             compatible; the attribution rests on the heading and on the
%             guard (v. 20, « Sophie Germain (autogr.) »), not on this pass.
%   v. 24–28  Two telegrams on printed Italian forms, received at Rome and
%             addressed « BALTHASAR BONCOMPAGNIE [sic] PALAIS PIOMBINO ROME »:
%             « RECU TELEGRAMME ACHAT AUTOGRAPHE GERMAIN » (Paris, 28 June
%             1884; v. 27–28) and « GERMAIN ACHETE 29 PART DEMAIN » (Paris,
%             30 June 1884; v. 24–26). Each is photographed with its address
%             flap down and up; v. 26 repeats v. 24 and v. 28 repeats v. 27,
%             noted only. They record the purchase of the Germain piece for
%             Boncompagni, which the slip on the guard at v. 20 also records.
%   v. 29     Blank opening. Not transcribed.
%   v. 30     Guard leaf for the next piece: lot « 335 », sale-catalogue
%             cutting no. 115 attributing the Latin verses to Constantin
%             Huygens, « Vendita Grimm. Romanzoff(?). Paris, mai 1886 ».
%   v. 31–32  Latin epigrams on Richelieu, a calligraphic seventeenth-century
%             hand, two leaves paginated 21 and 22: « In reditum Card.
%             Richelij », « In detractorem(?) Richelij mortui », « In illud
%             Borbonij », « In Richelij Bibliothecam », « In ejusdem obitum »
%             (headed « Audi alteram partem »), « In Regem », a distich after
%             an ampersand, and an untitled piece to Marie de Médicis. The
%             leaves name no author; the attribution is the cutting's.
%             Views 35 and 36 are second exposures of v. 31 and v. 32.
%   v. 33     Guard leaf for Mersenne: lot « 459 », cutting no. 98
%             (« 1° L. a. s. au Père Petau … 2° Pièce autographe, 1 p. 1/2
%             in-fol. »), « Vente Charavay, 6 février 1889 ». Its verso is the
%             blank left page. v. 37 repeats v. 33.
%   v. 34     Blank opening. Not transcribed.
%   v. 38–39L The « pièce autographe »: a folio in a small fast Latin cursive,
%             « Anno 1621 3 Januarij stilo. » — the declination of the
%             quadrangle star of the Little Bear and the elevation of the pole
%             at Paris, 48°47'35'', solved « Mechanicè per regulam et
%             circinum » and « Per Numeros », with a figure; its verso (v. 39
%             left) carries « Aliæ eiusd. anni observationes » (2–5 January
%             1621, polar star, pole 48°45'37½''; 23 and 25 December 1620,
%             equator altitude 41°18'20''/22'') and, on the folded lower half,
%             the address « Au Reverend pere Denis Petau de la Société de
%             Jesus Au College de Clermont ». Unsigned and undated beyond the
%             observations; « Anno 1621 » is its own.
%   v. 39R–40 The letter, on a small oblong leaf mounted alone: Mersenne to
%             Petau, in French, « Mon Reverend(?) pere … », sending
%             observations of 1621 on the height of the pole at Paris, asking
%             them back, and asking him to pass a word to « le R. pere
%             Grandamy(?) » on the longitude of Paris; signed « … Mar. Mers.
%             M. ». No place, no date on the leaf. A later hand below:
%             « Cette lettre est du P. Mersenne, minime ». v. 40 shows the
%             leaf turned back, its dorse bearing only « Au R. … pere Petau ».
%             Most of this very fast hand is left \ill{} or \uncertain{}.
%
% Numbers on the leaves. No thin pencil foliation of the library was found on
% any view, so no \folio{} is written. What is there, all in ink unless said:
% on each piece's leaves, top right, a running series 16 (v. 21), 17 (v. 22),
% 18 (v. 24–26, telegram), 19 (v. 27–28, telegram), 20 (guard v. 30),
% 21 (v. 31), 22 (v. 32), 23 (guard v. 33), 24 (v. 38), 25 or 26 (v. 39 R, the
% second figure looped), one per leaf of the album, versos unnumbered. It
% occupies the foliation's place but is ink and larger; it is recorded in
% notes only. A second, smaller series stands above it on the Huygens and
% Mersenne leaves: 547 (v. 31), 548 (v. 32), 319 (guard edge, v. 33/35),
% 320 (v. 38), 321 (v. 39 R) — and a centred 54 (v. 30), 21 (v. 33). On the
% telegrams, and in pencil on the Germain leaves, a fraction-like mark: 65
% over a bar with 6 (v. 21–22, 24–26) or 7 (v. 27–28) — apparently a lot and
% piece number; that meaning is an inference. The label on v. 21–22 reads
% « 16?9 », third figure ambiguous. « R.D. 9555 » stands on every Huygens and
% Mersenne leaf and on v. 21. Guard lot numbers: 335 (Huygens), 459
% (Mersenne); the sale notes carry « 64 » (v. 30) and « 19 » (v. 33).
%
% Hands and conventions. Germain's French keeps her spelling (connoître,
% differens, ennoncé, parceque, ceque, déduit for déduites); words written in
% a paler ink inside her lines are noted, not marked. The long s is
% transcribed s throughout. The Latin epigrams' hand writes final e and
% the second letter of æ as a small raised loop, final o like a round s,
% final t like a long-tailed s, and one looped capital for B, D, V and H; the
% readings follow the sense and are flagged where it does not settle them.
% In the Mersenne piece, the figure's Greek letters are set as Greek, the
% sun as $\odot$; the zodiac signs, for which LaTeX has no standard glyph,
% are set as their French names in brackets, \add{Bélier}, \add{Cancer},
% \add{Lion}, \add{Capricorne}, each note saying what the leaf draws; degrees, minutes and seconds keep the writer's
% points and primes.

\documentclass[11pt,a4paper]{article}
% The preamble shared by every transcription of the mathematicians' manuscripts in Gallica.
%
% It rests on one idea: **uncertainty compounds, it does not resolve.** A
% transcription that smooths over a doubtful word has destroyed the only thing
% it was made for — a reader must be able to tell, at every point, what was on
% the page and what was supplied. The apparatus macros below are therefore the
% critical apparatus, not decoration.
%
% The same commands are recognised by `scripts/render.mjs`, which produces the
% site's reading view, and by `scripts/tei.mjs`. A macro added here must be
% added there too, or rendering fails loudly — which is the wanted behaviour.
%
% Comments here are English, like the rest of the repository. The documents
% this preamble sets are French, and that is deliberate: see the skills.

\ProvidesPackage{germain}[2026/09/15 Transcriptions of mathematicians' manuscripts in Gallica]

\usepackage{amsmath,amssymb,amsthm}
\usepackage{mathrsfs}
% Kept from the parent preamble so the renderer's diagram support stays
% available; these manuscripts draw no commutative diagrams, and nothing here uses it.
\usepackage{tikz-cd}
\usepackage[english,french]{babel}
\usepackage{fontspec}

% --- Greek ------------------------------------------------------------------
% The text face stays fontspec's default, Latin Modern, which has no Greek:
% XeTeX drops every glyph it lacks with a line in the log and nothing on the
% page, and the Greek words the manuscripts quote vanished from the PDFs. So
% the Greek and Greek Extended blocks get a XeTeX character class of their own,
% and entering or leaving a run of it switches the family to CMU Serif — the
% same Computer Modern design, with polytonic Greek — and back. Only the
% family changes: a Greek word in \textit or \textbf keeps its shape and series.
% The transitions to and from every other class carry the tokens that class
% has towards an ordinary letter, so French spacing before « ; » or inside
% « » still holds after a Greek word. No grouping, so no brace or \endgroup
% can come out unbalanced; maths is left alone, since XeTeX inserts nothing
% there. Kept identical in `exercices.sty`.
\makeatletter
\newfontfamily\germainGreekfont{cmun}[
  Extension=.otf, NFSSFamily=germaingreek,
  UprightFont=*rm, ItalicFont=*ti, SlantedFont=*sl,
  BoldFont=*bx, BoldItalicFont=*bi, BoldSlantedFont=*bl]
\newXeTeXintercharclass\germain@greek
\def\germain@greekrange#1#2{%
  \@tempcnta=#1\relax
  \loop
    \XeTeXcharclass\@tempcnta=\germain@greek
    \ifnum\@tempcnta<#2\advance\@tempcnta\@ne\repeat}
\germain@greekrange{"0370}{"03FF}
\germain@greekrange{"1F00}{"1FFF}
\def\germain@greekprev{\rmdefault}
\def\germain@greekin{%
  \edef\germain@greekprev{\f@family}\fontfamily{germaingreek}\selectfont}
\def\germain@greekout{\fontfamily{\germain@greekprev}\selectfont}
\def\germain@greekjoin{%
  \XeTeXinterchartokenstate=\@ne
  \@tempcnta=\z@
  \loop
    \ifnum\@tempcnta=\germain@greek\else
      \XeTeXinterchartoks\germain@greek\@tempcnta=\expandafter{%
        \expandafter\germain@greekout\the\XeTeXinterchartoks\z@\@tempcnta}%
      \XeTeXinterchartoks\@tempcnta\germain@greek=\expandafter{%
        \the\XeTeXinterchartoks\@tempcnta\z@\germain@greekin}%
    \fi
    \ifnum\@tempcnta<4095\advance\@tempcnta\@ne\repeat}
% babel-french sets its own classes, and saves and restores the interchar
% state, each time a language is selected: join again after it does.
\addto\extrasfrench{\germain@greekjoin}
\addto\extrasenglish{\germain@greekjoin}
\AtBeginDocument{\germain@greekjoin}
\makeatother

\usepackage{geometry}
\geometry{a4paper,margin=25mm,marginparwidth=28mm}
\usepackage{marginnote}
\usepackage{xcolor}
\usepackage[normalem]{ulem}
\setlength{\emergencystretch}{3em}
\usepackage{hyperref}
\hypersetup{colorlinks=true,linkcolor=black,urlcolor=black,pdfborder={0 0 0}}

% --- Batch metadata ---------------------------------------------------------
% Taken as they stand by the rendering script, which shows them at the head of
% the reading view. They fix what the file claims to cover. The macro names are
% the parent project's, so that the scripts stay shared; \folder here means the
% volume's slug — `fr-9115` — and \pages counts Gallica views.

\newcommand{\folder}[1]{\gdef\grothFolder{#1}}
\newcommand{\batch}[1]{\gdef\grothBatch{#1}}
\newcommand{\pages}[2]{\gdef\grothFirst{#1}\gdef\grothLast{#2}}
\newcommand{\foldertitle}[1]{\gdef\grothFoldertitle{#1}}

% The holder's own shelfmark, printed as they write it: « Français 9115 ».
\newcommand{\shelfmark}[1]{\gdef\grothShelfmark{#1}}

% The dating, copied verbatim from the holder's catalogue — never inferred
% here. « XIXe siècle » is the BnF's; a pass that dates a leaf from its content
% says so in a \note{}, as its own inference.
\newcommand{\dating}[1]{\gdef\grothDating{#1}}

% The status notice, printed in the title block and repeated as a diagonal
% watermark on every page of the PDF. The manuscripts are in the public
% domain; what this line declares is not a right but a quality — an
% unverified first pass by a machine. The skills write the line; a file
% without it gets no watermark, so leaving it out is visible in the output.
\newcommand{\watermark}[1]{\gdef\grothWatermark{#1}}

\gdef\grothFolder{?}\gdef\grothBatch{}\gdef\grothFirst{?}\gdef\grothLast{?}%
\gdef\grothFoldertitle{}\gdef\grothShelfmark{}\gdef\grothDating{}\gdef\grothWatermark{}

% --- The résumé -------------------------------------------------------------
\newenvironment{resume}
  {\small\setlength{\parskip}{0.45em}}
  {\par\medskip\hrule\medskip}

% --- Keywords ---------------------------------------------------------------
% In English, closing the modernised reading's résumé: the modern vocabulary
% under which to search for what the leaves construct. The manifest extracts
% them to tag the volume — this line is the single source of the tags.
\newcommand{\keywords}[1]{\par\smallskip\noindent
  {\footnotesize\textsc{Keywords} --- \textit{#1}}\par}

% --- The critical apparatus -------------------------------------------------

% The Gallica view number, in the margin. This is the anchor that lets a
% disagreement over a formula be settled by looking at *one* image rather than
% a volume of 750. The site uses it to turn the facsimile as the transcription
% is scrolled. A view is one scanned image: usually one side of a leaf.
\newcommand{\page}[1]{%
  \marginnote{\footnotesize\color{gray!70}\textsf{v.\,#1}}%
  \label{page:#1}\ignorespaces}

% The BnF's pencilled foliation, where the leaf carries it and a pass can read
% it: « 198r ». This is what the literature cites — « ff. 198r–208v » — and
% the one line that turns a citation into a view. Recorded, never inferred:
% a folio a pass cannot read is not written.
\newcommand{\folio}[1]{%
  \marginnote{\footnotesize\color{gray!50}\textsf{f.\,#1}}\ignorespaces}

% The modernised reading's coarser counterpart to \page{}: which views a
% section is drawn from.
\newcommand{\pagerange}[2]{%
  \marginnote{\footnotesize\color{gray!70}\textsf{v.\,#1--#2}}%
  \ignorespaces}

% Illegible: never guessed.
\newcommand{\ill}{\textcolor{red!60!black}{[\,\ldots\,]}}

% A doubtful reading: the word is given, and flagged as uncertain.
\newcommand{\uncertain}[1]{\uline{#1}}

% An editorial addition — a missing letter, an implied word.
\newcommand{\add}[1]{\textcolor{blue!50!black}{[#1]}}

% Struck out by the writer. What was crossed out is part of the document.
\newcommand{\struck}[1]{\sout{#1}}

% The transcriber's note, on the state of the page or the mathematics.
\newcommand{\note}[1]{\par\smallskip{\small\color{gray!60!black}\textit{[#1]}}\par\smallskip}

% A marginal note of hers, distinct from the body of the text.
\newcommand{\marginal}[1]{\par\smallskip\noindent\hspace{1em}%
  {\small\color{gray!60!black}#1}\par\smallskip}

% --- Title ------------------------------------------------------------------
% The title block is French, like the documents themselves.

\AddToHook{shipout/background}{%
  \ifx\grothWatermark\empty\else
    \begin{tikzpicture}[remember picture, overlay]
      \node[rotate=52, scale=3.4, align=center, text=gray!18,
            font=\sffamily\bfseries]
        at (current page.center) {\grothWatermark};
    \end{tikzpicture}%
  \fi}

\AtBeginDocument{%
  \begin{center}
    {\large\bfseries Manuscrits de mathématiciens (Gallica) ---
      \ifx\grothShelfmark\empty\grothFolder\else\grothShelfmark\fi}\\[2pt]
    {\itshape\grothFoldertitle}\\[4pt]
    {\small vues \grothFirst--\grothLast
      \ifx\grothBatch\empty\else\ (lot \grothBatch)\fi}\\[2pt]
    \ifx\grothDating\empty\else
      {\small\color{gray!60!black}Datation du catalogue : \grothDating}\\[2pt]
    \fi
    \ifx\grothWatermark\empty\else
      {\small\bfseries\color{red!45!gray}\grothWatermark}\\[2pt]
    \fi
    {\footnotesize\color{gray!60!black}Transcription établie d'après les images IIIF de Gallica.
     Source gallica.bnf.fr / Bibliothèque nationale de France.\\
     Les lectures incertaines sont soulignées, les illisibles notées [\,\ldots\,],
     les ajouts entre crochets ; \textsf{v.} numérote les vues de Gallica,
     \textsf{f.} la foliotation au crayon de la BnF.}
  \end{center}
  \vspace{1em}}

\endinput


\folder{naf-9544}
\shelfmark{NAF 9544}
\batch{2}
\pages{21}{40}
\dating{1601-1700}
\watermark{Lecture automatique\\non vérifiée}
\foldertitle{Recueil de lettres autographes de D'ALEMBERT, l'abbé BOSSUT, Sophie GERMAIN, C. HUYGHENS et le P. MERSENNE.}

\begin{document}

\page{21}
\note{Page de droite de l'ouverture ; la page de gauche est un onglet blanc. Un petit feuillet est monté sur l'onglet. En tête, entre crochets, une ligne d'une plume plus épaisse. En haut à droite, à l'encre, « 16 » ; au-dessous, au crayon et très pâle, un chiffre (6 ou 65) posé sur un trait horizontal, et au bord droit une autre figure au crayon qui n'a pas été lue. Sur une étiquette collée à droite du feuillet, à l'encre, « 16\uncertain{5}9 », le troisième chiffre étant d'une forme ambiguë. Dans la marge de gauche, au bas du feuillet, d'une autre main : « R.D. 9555 ». Le cachet rond de la bibliothèque est apposé au milieu du texte. Aucune de ces marques n'est donnée pour la foliation de la bibliothèque.}

[Ce qui suit s'est trouvé dans les papiers laissés par M\textsuperscript{lle} Germain].

L'idée fondamentale d'une science n'est pour celui qui en commence l'étude que la déterminante du sujet de cette science ; à mesure qu'il avance, il apprend a connoître les differens points de vues sous lesquels l'objet de la science peut etre considéré et lorsqu'il est parvenu a la connoissance parfaite du sujet l'idée fondamentale est pour lui la science toute entière, parceque cette idée est alors agrandie de toutes celles qu'il a acquis depuis, et de leurs rapports.

Mais a-t-on dit ? si la science toute entière n'étoit pas contenue dans l'idée fondamentale comment pourroit \struck{on} l'y voir
\note{le mot après « pourroit » est chargé d'encre au point d'être effacé ; la lecture « on » est celle du sens, et il n'est pas sûr qu'il soit biffé plutôt que repassé.}

Je reponds ; La proposition fondamentale d'une science, bien que contenue dans des termes exacts et constans ne renferme pas pour tous ceux qui l'entendent une idée identiquement la même mais seulement des idées qui conviennent dans leurs traits caractéristiques de sorte que pour tous les esprits le sujet est egalement déterminé mais inégalement connu et que cette proposition qui n'a d'abord été \add{\uncertain{à l'entrée de la carrière}} qu'une idée caractéristique devient lorsqu'on l'a parcourue, une veritable description du sujet et remplit au plus haut degré la fonction de définition par rapport a la science dont elle détermine l'objet.
\note{« Je » en tête de ligne est à peine formé. Les mots donnés comme ajout sont écrits au-dessus de la ligne, entre « été » et « qu'une », d'une plume plus fine : le premier est un signe bref qui peut être « à », les derniers se lisent mal. « parrapport » est écrit en un seul mot.}


\page{22}
\note{Le feuillet de la vue 21 est déplié : la page de gauche en est le verso, la page de droite le recto d'un second feuillet, qui porte en haut à droite, à l'encre, « 17 », et au crayon pâle, à sa gauche, un nombre (65) posé sur un trait, avec un chiffre au-dessous. L'étiquette « 16\uncertain{5}9 » de la vue 21 dépasse à droite. Le cachet « BIBLIOTHÈQUE NATIONALE R.F. MSS » est apposé au milieu de la page de droite. Plusieurs mots, sur les deux pages, sont d'une encre nettement plus pâle que le reste de la ligne (« la réunion », « se lie », « de l'idée » à gauche ; « comment », « leur demonstration », « première », « repondre », « toutes » à droite), comme s'ils avaient été écrits dans des blancs laissés d'abord.}

En effet le souvenir des propositions qui ont été déduit de la proposition fondamentale et dont la réunion forme l'ensemble de la science se lie pour celui qui la possede a l'ennoncé de l'idée fondamentale de sorte qu'il voit d'un seul coup d'œil le sujet tout entier.

On a confondu l'idée fondamentale avec le sujet de la science c'est ainsi que l'on a pu dire ; \emph{l'inconnu est dans le connu ; lorsque l'on apprend on va du même au même ; une science ne renferme qu'une seule idée} :
\note{les trois propositions citées sont soulignées. Un petit signe, peut-être un renvoi, est écrit au-dessus du « même » qui ouvre la ligne. « fondamentale » est tracé ici d'une main pressée, le milieu du mot en partie redoublé.}

Rien n'est plus absurde que de \struck{supposer} \add{dire} que l'on puisse avoir une idée sans la connoître et c'est cependant ceque cette doctrine suppose.
\note{« dire » est écrit au-dessus de « supposer » biffé.}

On a été jusqu'à \struck{dire} \add{prétendre} que dans une suite de déductions vers la seconde ou troisième il ne restoit plus d'idée, mais seulement des signes d'idées, puis ensuite des signes de signes, \struck{puis} et qu'un raisonnement bien fait n'étoit qu'une operation mécanique exécutée sur des signes qu'il falloit pour revenir aux idées \struck{arriver} \uncertain{jus} retourner jusqu'à l'idée fondamentale.
\note{« prétendre » est écrit au-dessus de « dire » biffé. La page de gauche finit sur « arriver », barré d'un trait, suivi de deux ou trois lettres que le même trait semble traverser ; la phrase reprend en tête de la page de droite par « retourner jusqu'à ».}

Et a-t-on dit, si vous croyez avoir des idées dans une longue suite de deductions comment ne vous contentez vous sur leur demonstration qu'en revenant a l'idée première.

Je crois que l'on peut repondre ; les propositions sont essentiellement des idées toutes les fois qu'elles ont un sens intelligibles, si dans une suite de déductions on sent le besoin de retourner en arrière c'est uniquement pour s'assurer qu'elles conviennent au sujet auquel on les attribue c'est enfin pour s'assurer de l'identité du sujet ennoncé avec le sujet qualifié.

La certitude mathématique tient a l'identité du sujet et la langue des calculs est bien faite parceque en raisonnant suivant les regles qu'elle prescrit il est impossible d'attribuer au sujet de la science des propriétés qui ne soient pas contenues dans l'idée fondamentale de sorte que les questions que l'on examine ne se trouvent
\note{la page s'arrête en pleine phrase, sur un trait en forme de croix double après « trouvent ». La phrase s'achève en tête du verso, vue 23.}


\page{23}
\note{Verso du feuillet paginé 17 (vue 22). Seules les deux premières lignes sont écrites ; le reste de la page n'est que le texte du recto vu par transparence, à l'envers. L'étiquette « 16\uncertain{5}9 » est visible à droite, sur la page de garde.}

jamais compliquées de considerations qui lui soient étrangères
\note{fin de la phrase commencée au bas de la vue 22 ; le fragment s'arrête ici, sans signature ni date.}


\page{24}
\note{Feuillet de l'album photographié en hauteur, la moitié supérieure blanche. Dans la moitié inférieure est monté un télégramme reçu à Rome, sur une formule imprimée de l'administration italienne (« Indicazioni eventuali abbreviate (Mod. 30) », « Ufficio Telegrafico di », « Ricevuto il … ore … Pel circuito N.° … Ricevente »), le rabat d'adresse relevé. Dans la case « Ufficio Telegrafico di », un cachet dont ne se lit que l'initiale R. Le texte de la dépêche est imprimé sur des bandes de papier collées ; les mentions de réception sont à la main. Cachet rond « BIBLIOTHÈQUE NATIONALE MSS » sur la bande d'en-tête. Au coin inférieur droit, écrits dans le sens de la hauteur, « 18 » à l'encre et un nombre « 65 » accompagné d'un 6, séparés par un trait.}

Ricevuto il 30/6 1884 ore 19 \uncertain{1/4} — Pel circuito N.° 5 — Ricevente \uncertain{Savoz}
\note{mentions manuscrites portées dans les cases de la formule ; l'heure est suivie d'un signe bref qui peut être une fraction.}

ROME PARIS 42442 10 30 4/35 SR .+

.+ GERMAIN ACHETE 29 PART DEMAIN .-
\note{les deux lignes en capitales sont des bandes imprimées par le télégraphe et collées sur la formule. Selon l'avis imprimé de la formule, les nombres qui suivent le lieu d'origine donnent le numéro de la dépêche, le nombre de mots, puis la date, l'heure et les minutes du dépôt : dépêche 42442, dix mots, déposée le 30 à 4 h 35.}

\page{25}
\note{Le même télégramme, rabat d'adresse rabattu : la bande collée sur le rabat porte le nom du destinataire. Au-dessus, la mention imprimée « Il porto è gratuito » et le numéro de remise, à l'encre. Le timbre de fermeture est collé au bas du rabat. Au coin inférieur droit, les mêmes « 18 » et « 65 » que sur la vue 24.}

\uncertain{571}
\note{numéro manuscrit porté devant « N.° … di recapito ».}

BALTHASAR BONCOMPAGNIE PALAIS PIOMBINO

ROME .
\note{« BONCOMPAGNIE » est ainsi imprimé sur la bande.}

\page{26}
\note{Seconde photographie du télégramme des vues 24 et 25, rabat relevé, identique à la vue 24 : mêmes mentions, mêmes bandes, même « 18 » au coin. Rien n'y est transcrit de nouveau.}

\page{27}
\note{Feuillet suivant de l'album, en hauteur, la moitié supérieure blanche : un second télégramme, sur la même formule, rabat d'adresse rabattu. Numéro de remise à l'encre sur le rabat ; dans la case « Ufficio Telegrafico di », un cachet oblique « \ill{}MA CENTRALE ». Cachet rond « BIBLIOTHÈQUE NATIONALE MSS » sous la bande d'en-tête. Au coin inférieur droit, écrits dans le sens de la hauteur, « 19 » à l'encre et un nombre « 65 » accompagné d'un 7.}

\uncertain{657}

BALTHASAR BONCOMPAGNI PALAIS

PIOMBINO ROMA =

Ricevuto il 28/6 188\uncertain{4} ore 21.20 — Pel circuito N.° 1 — Ricevente \ill{}
\note{le dernier chiffre de l'année, ajouté à la main après le « 188 » imprimé, est mal formé. Le nom du receveur, bref, n'a pas été lu.}

ROME PARIS NO 32042 10 28/6/84 6/30 SR

RECU TELEGRAMME ACHAT AUTOGRAPHE GERMAIN =
\note{bandes imprimées par le télégraphe et collées. D'après les dates qu'elles portent, cette dépêche du 28 juin 1884 précède celle des vues 24–26, datée du 30.}

\page{28}
\note{Le même télégramme que la vue 27, rabat relevé : le timbre de fermeture est collé sur le rabat. Les mentions et les bandes sont celles de la vue 27 ; rien de nouveau n'y est transcrit.}


\page{30}
\note{Feuillet de garde de l'album, sans texte ancien : il annonce la pièce qui suit. En haut au milieu, « 54 » à l'encre ; en haut à droite, « 20 », et tout contre le bord un chiffre plus petit, rogné. Au-dessous, un grand « 335 » à l'encre. Collée au milieu, une coupure imprimée d'un catalogue de vente, numéro 115 : « Huyghens (Constantin), seigneur de Zuylichem, célèbre poète et homme d'État hollandais, correspondant de Descartes, de Balzac et de Pierre Corneille, né en 1596, mort en 1687. Pièce de vers aut., en latin, 3 p. in-fol. Epigrammes sur le retour de Richelieu, sur sa mort, sa bibliothèque ; d'autres sont dirigées contre Louis XIII, qui doit craindre le triomphe des Espagnols après la mort de Richelieu ; la dernière est adressée à Marie de Médicis. » À gauche de la coupure, « R.D. 9555 ». Au-dessous, à l'encre : « Vendita Grimm. \uncertain{Romanzoff}. Paris, mai 1886 » — mot italien en tête d'une ligne française ; puis « 64 ». La page de gauche de l'ouverture est blanche.}


\page{31}
\note{Page de gauche : blanche. Page de droite : un feuillet de vers latins, d'une main du XVII\textsuperscript{e} siècle, calligraphiée, à grands paraphes sous les lignes. En haut à droite, « 21 » à l'encre, et au-dessus, plus petit, « 547 ». Au bas de la page, d'une main moderne, « R.D. 9555 », souligné d'un trait ondulé. Cachet rond de la Bibliothèque nationale au milieu du second poème. La coupure de catalogue collée sur la garde qui précède (vue 30) attribue ces vers à Constantin Huygens.}

\emph{In reditum Card. Richelij.}

Pallet atrâ spe lapsus Iber : quo plurima mundi

Machina versatur Cardo \uncertain{prodigio} agit.

Pristina Richelijs redierunt robora nervis.

Cessit Atlanteis pestis adusta toris.

Solve duci votum reduci, quo vincere pro te,

Quo te complecti, Gallia, pergat \uncertain{habet}.
\note{Cette main écrit le e final comme une petite boucle surélevée (« Solve », « vincere », « te », « spe ») et le o final comme un s arrondi (« Cardo », « salso ») ; la lecture se règle sur le sens. Le mot après « Cardo » se présente comme « ouodigie » ; « prodigio » est offert avec doute. Le dernier mot du distique, suivi de deux points, se présente comme « Labet » : son initiale est la capitale bouclée que la même main emploie pour le H de « hactenus » à la vue 32, d'où la lecture « habet », offerte avec doute ; le même mot clôt un vers de la vue 32.}

\emph{In \uncertain{detractorem} Richelij mortui.}

Quæ Cardinali turba supplicem vivo

Subjicit et mentem et manum, nec in terris

Præsentius prostrata numen agnovit,

\uncertain{Ruit} in sepultum scommatis, jocis, probris,

Et quidquid ore possit expui salso

In immerentis impotens vomit manes.

Levis popelle, ingrata plebs, inhumanum

Serpentis instar ; Patriæ simel natus,

Atlas, Atrides, Cardo Regis et Regni,

Non debuit denascier ; tibi verò

Nec debuit defunctus esse, nec vivus.
\note{Le t final de cette main ressemble à un s à longue queue (« Subjicit », « agnovit », « possit », « vomit », « debuit ») ; il est lu t partout où « debuit » le fixe. Le premier mot du quatrième vers a pour initiale une grande capitale bouclée, Q ou R : « Ruit » est offert par le sens. « simel » est écrit ainsi. Æ est rendu par a suivi de la petite boucle surélevée (« Quæ », « Præsentius », « Patriæ »). Le titre du second poème se termine par cette même boucle, suivie de deux points.}


\page{32}
\note{Ouverture : à gauche, le verso du feuillet de la vue 31 ; à droite, un second feuillet de la même main, qui porte en haut à droite « 22 » à l'encre et, plus petit, « 548 ». Cachet rond de la Bibliothèque nationale au milieu de la page de droite. Chaque pièce se termine par un trait oblique. Sur la page de gauche, le texte du recto transparaît.}

\emph{In illud Borbonij}

\emph{Ext\textsuperscript{o} subortis esca vermibus \uncertain{late}.}
\note{la seconde ligne du titre, soulignée, cite le vers de Bourbon auquel l'épigramme répond. « Ext » est suivi de la petite boucle surélevée de cette main, abréviation dont la résolution n'est pas donnée ; le dernier mot est d'une lecture douteuse.}

Vellitur indignè prostrato barba Leoni ;

Richelium plebis improba dente ferit ;

\uncertain{Vincit}, sed extinctum : soboles, quæ supplice vivum

Nuper adoravit, mordicus ore \uncertain{patrem}.

Ilicet, sic tantum non falleris impie, natis

Ingratis de se vermibus esca jacet.
\note{le V initial du troisième vers est de la forme bouclée qui sert aussi pour B et D dans cette main. Le dernier mot du quatrième vers se présente comme « patrim ».}

\emph{In Richelij Bibliothecam.}

Mille libros, et mille libros, et mille librorum

Millia magnifici publica Richelij

Cura, bono patriæ, capsis ingesta superbis

Nominis æterni marmora viva locat.

Quæque voluminibus moles augusta \uncertain{coemptis}

Cresceret in titulos auctior usque suos,

Mille libros et mille libros, et mille librorum

Jussit adimpleri millia mille libris.

Desine posteritas \ill{} \struck{\ill{}}, quo se agmine facto

Duplicet in multam massa voluta nivem.

Materiam liquit vestris qui scribitis, æquam

Viribus ; omnis ubi Gallia sudet \uncertain{habet}.

Scribite Richelium, partos pro Rege triumphos

Scribite et æternis gesta legenda typis.

Mille libros impletura est quæ postuma famam

Unius capiet Bibliotheca viri.
\note{« coemptis » se présente comme « coimptis ». Après « posteritas », un mot de cinq ou six lettres, un a surmonté d'un petit rond, puis « rroe », n'a pas été lu ; il est suivi d'un mot bref barré d'un trait horizontal. Le dernier mot du vers « omnis ubi Gallia sudet … » a pour initiale la capitale bouclée que cette main emploie aussi pour le H de « Hactenus » (page de droite) : « habet » est offert avec doute, comme à la vue 31.}

\note{En tête de la page de droite, à gauche, en plus petit et de la même main : « Audi alteram partem. »}

\emph{In ejusdem obitum.}

Sic decet : Armandus dominæ comes ivit \uncertain{Hetruscæ}.

Galle, quid insano degener ore gemis ?

Quod sis esse velis. tandem sub Rege monarcha es.

Successit matri Filius et famulo.

\emph{In Regem.}

Qui scipionem perdidit, cave ne nunc

Plus in die quàm septies cadat Justus.

\&

Hactenus Armandum suspexit non malè Regem :

Nunc exarmandum non malè sperat Iber.
\note{le signe entre les deux pièces est une esperluette isolée, au milieu de la ligne.}

Maria, mater impotentis gratiæ,

Quocumque manes inter \uncertain{ut} parta simul

\uncertain{Quiete} vivæ denegata consides,

Cave Cardinalem, si sapis : suspectus est,

Qui tam repente mortuam voluit sequi,

Quam persecutus hactenus vivam fuit.
\note{pièce sans titre. « Quiete » se présente comme « Quieti » suivi de la petite boucle surélevée ; le mot entre « inter » et « parta » peut être « ut » ou « et ». Une tache d'encre couvre le début de « consides ». « hactenus » est écrit avec la même capitale bouclée que le « habet » supposé de la page de gauche.}


\page{33}
\note{Page de gauche : le verso du feuillet paginé 22, blanc ; le texte du recto y transparaît à l'envers. Page de droite : feuillet de garde qui annonce la pièce suivante. En haut à droite « 23 » à l'encre, et au bord extrême un nombre plus petit, « 319 » ; en haut au milieu, « 21 ». Au-dessous, un grand « 459 » souligné d'un long trait. Collée, une coupure imprimée d'un catalogue de vente, numéro 98 : « Mersenne (le Père Marin), célèbre géomètre, physicien et philosophe, ami et défenseur de Descartes, né dans la Sarthe, 1588, m. 1648. 1° L. a. s. au Père Petau, 1 p. in-4 oblong. Belle lettre où il lui annonce la résolution d'un problème relatif à la hauteur du pôle au-dessus de l'horizon de Paris. 2° Pièce autographe, 1 p. 1/2 in-fol. C'est la résolution du problème annoncée dans la lettre précédente. » À gauche de la coupure, « R.D. 9555 ». Au-dessous, à l'encre : « Vente Charavay, 6 février 1889 », puis « 19 ».}


\page{35}
\note{Seconde photographie de la page de la vue 31 (même feuillet paginé 21, mêmes vers, même « R.D. 9555 ») ; le cadrage montre en plus, au bord supérieur droit, le petit nombre « 319 » de la garde voisine. Rien n'y est transcrit de nouveau.}


\page{36}
\note{Seconde photographie de l'ouverture de la vue 32 (verso du feuillet 21 et feuillet paginé 22) : même texte. Rien n'y est transcrit de nouveau.}


\page{37}
\note{Seconde photographie de l'ouverture de la vue 33 (garde paginée 23 et coupure « Mersenne »). Rien n'y est transcrit de nouveau.}


\page{38}
\note{Page de gauche : blanche. Page de droite : un feuillet in-folio écrit d'une petite cursive rapide du XVII\textsuperscript{e} siècle, en latin, sous une figure. La coupure de la garde (vue 33) le décrit comme la « Pièce autographe, 1 p. 1/2 in-fol. » de Mersenne, résolution du problème annoncé dans la lettre au P. Petau. En haut à droite, « 24 » à l'encre, et au-dessus, plus petit, « 320 ». Au bas, d'une main moderne, « R.D. 9555 » souligné d'un trait ondulé. Cachet rond de la Bibliothèque nationale au milieu du texte. Le coin supérieur droit est taché et rongé.}

\note{Figure, en haut de la page : un grand cercle, le méridien, de diamètre horizontal $\alpha\beta$ (l'équinoxial) et de centre I ; les diamètres $\varepsilon\zeta$ (le zodiaque), $\gamma\delta$ (l'axe du monde), $\eta\theta$, et un diamètre de $\rho$ à s. Au sommet, près de $\gamma$, deux demi-cercles qui s'entrecroisent au-dessus du cercle, et les points o, $\gamma$, A, E, F, $\xi$, $\mu$ ; plus bas, sur la verticale issue de A, les points D, C, B, B sur $\alpha\beta$ ; la perpendiculaire abaissée de $\kappa$ tombe en $\lambda$ sur $\alpha\beta$. Des cordes horizontales joignent D et C au cercle. Seules les lettres sont données ici ; le tracé n'est pas reproduit.}

Anno 1621 3 Januarij stilo. Lucida quadranguli Cynosuræ visa supra horizontem in meridiano sub polo alta partium 34.35$'$. Eodem tempore longitudo stellæ in 7.34$'$.20$''$ \add{Lion}, et ideo ab æquinoctio verno 127.34$'$.20$''$. Latitudo stellæ 72.51$'$.30$''$. Quæritur iam declinatio stellæ, quam hac figura \uncertain{exhibeo}.
\note{« stilo » est suivi d'un point, sans qualificatif. Le signe placé après la longitude est celui du Lion, rendu \add{Lion} ; le même signe revient plus bas, avec celui du Cancer.}

$\alpha\gamma\beta\delta$ Meridianus est circulus. $\alpha\beta$ æquinoctialis circulus. $\gamma\delta$ axis mundi, et ideo $\gamma$ polus arcticus, et $\delta$ polus antarcticus. $\varepsilon\zeta$ Zodiacus. $\eta$ polus Zodiaci Borealis. $\alpha\varepsilon$, sive $\beta\zeta$ maxima solis declinatio secundum \uncertain{Tychonem} partium 23.31$'$.30$''$. $\gamma\eta$ sive $\delta\theta$ eadem declinatio polorum Zodiaci $\eta\theta$ a polis mundi $\gamma\delta$ 23.31$'$.30$''$. $\varepsilon$A sive $\zeta\kappa$ arcus latitudinis stellæ utrinque a Zodiaco $\varepsilon\zeta$ supputatus 72.51$'$.30$''$. A$\kappa$ parallelus latitudinis stellæ. A$\mu\kappa$ parallelus circulus latitudinis stellæ. A$\mu$ longitudo stellæ super dicto parallelo supputata, vel a puncto æquinoctiali \add{Bélier} sive solstitiali A supputata partium 127.34$'$.20$''$. Vel a dicto solstitiali puncto \add{Lion} 7.34.20 ; vel 37.34$'$.20$''$ a \add{Cancer}.
\note{« maxima » est précédé d'un trait en boucle, peut-être un mot biffé. Le nom après « secundum » se présente comme « Tichonem » ou « Tirhonem » ; « Tychonem » est offert avec doute. Le signe qui suit « æquinoctiali » a la forme d'un y et est rendu \add{Bélier} ; celui qui clôt la phrase a la forme d'un 69 couché et est rendu \add{Cancer}. La main écrit la lettre A de la figure comme un $\lambda$ capital ; le $\lambda$ minuscule de la figure est distinct. La phrase « sive solstitiali A supputata partium 127.34.20 » est transcrite telle qu'elle se lit, bien que 127°34'20'' soit compté depuis l'équinoxe.}

Quæritur arcus $\alpha$o vel $\beta\xi$ \uncertain{uterque} arcus declinationis stellæ ab æquinoctiali $\alpha\beta$ supputata.

\emph{Mechanicè \uncertain{per} regulam et circinum}
\note{titre centré. Entre « Mechanicè » et « regulam », un p suivi d'un point et d'un signe qui peut être « et » : la lecture n'est pas sûre.}

Singulis iuxta dimensionem partibus ordinatis, iam a puncto longitudinis $\mu$ \uncertain{descendat} perpendicularis $\mu$F super parallelum latitudinis punctum F erit verus locus stellæ in cœlo. A puncto F \uncertain{excitetur} æquinoctiali $\alpha\beta$ parallela OF$\xi$ utrinque producta donec meridianum secet in O et $\xi$. Tunc arcus $\alpha$O, sive $\beta\xi$ repræsentabit \uncertain{instar} pro tempore stellæ declinationem 75.47.20. Arcus $\xi\gamma$ complementum declin. 14.12$'$.35$''$. Ideoque si a puncto $\xi$ infra polum arcticum $\gamma$ sumatur arcus $\xi\rho$ æqualis datæ altitudini in meridiano 34.35$'$, restabit $\gamma\rho$ arcus elevationis poli arctici supra locum observationis 48.47$'$.35$''$.
\note{90° moins 75°47'20'' donnent 14°12'40'', non 14°12'35'' ; la somme 48°47'35'' qui suit est celle de 14°12'35'' et de 34°35'. Transcrit tel qu'écrit. Le mot après « repræsentabit » commence par « ins » et n'est pas sûr.}

\emph{Per Numeros.}

Primo iungatur $\alpha\varepsilon$ max\textsuperscript{a} decl. solis cum $\varepsilon$A Latitudo stellæ, fiet arcus $\alpha$A 96.23. Ideo A$\beta$ 83.37 sinus rectus erit partium 99380, scilicet recta AB. 2\textsuperscript{o} auferatur $\zeta\beta$ ab illo, scilicet $\zeta\beta$ a $\zeta\kappa$, superabit $\kappa\beta$ 49.20, cuius sinus rectus $\kappa\lambda$ part. 75851 ablatus a superiori AB relinquit AC par. 23529, cuius medietas dat AD par. 11764. Tertio sinus versus arcus A$\mu$ longitudinis dempto sinu complementi eius e radio 100000, \uncertain{is erit} 20741 videlicet recta AF. Jam in triangulo rectangulo A\uncertain{D}G hæc est proportio ut AG ad AD sic AF ad AE. Est autem AG sinus totus 100000 scilicet radius circuli A$\mu\kappa$, ergo ut AG 100000 ad AD 11764, sic AF 20741 ad AE 2440. quo dempto ex AB 99380 restat EB 96940 sinus rectus declinationis stellæ, videlicet arcus $\alpha$o vel $\beta\xi$ 75.47.20. Ablata declinatione stellæ proposita 75.47.20 a toto quadrante 90 remanebit o$\gamma$ vel $\xi\gamma$ complementum declinat. par. 14.12$'$.35$''$ cui addita stellæ alt. merid. sub polo $\gamma$ partium 34.35 videlicet arcus $\xi\rho$, erit totus $\gamma\rho$ arcus elevationis poli supra locum \struck{48} id est Lutetiæ 48.47.35. Ergo s$\rho$ repræsentabit horizontem.
\note{« Latitudo » est au nominatif, comme écrit. Dans « A D G », la seconde lettre est repassée et peut être un B. Avant « id est Lutetiæ », un « 48 » est biffé d'un trait. Le dernier mot se présente comme « horizontem » écrit vite, l'h initial en longue boucle. Les nombres s'accordent entre eux : 99380 − 75851 = 23529, dont la moitié est 11764 ; 11764 × 20741 / 100000 donne 2440 ; 99380 − 2440 = 96940.}


\page{39}
\note{Ouverture. À gauche, le verso du feuillet paginé 24 (vue 38), écrit de la même petite cursive latine : une moitié supérieure d'observations, puis, dans la moitié inférieure, l'adresse, le feuillet ayant été plié en lettre. Le texte du recto transparaît dans les blancs. À droite, sur un onglet, un petit feuillet oblong monté seul : une lettre en français, qui porte en haut à droite « 321 » et, au-dessous, « 2\uncertain{5} » à l'encre, le second chiffre bouclé pouvant se lire 5 ou 6. À gauche de la signature, « R.D. 9555 » souligné d'un trait ondulé ; au-dessous, d'une autre main, une attribution. Le cachet rond de la Bibliothèque nationale est apposé au milieu de la lettre, sur plusieurs mots.}

\emph{Aliæ eiusd. anni observationes}
\note{titre centré, en tête de la page de gauche.}

Die 2. Januarij h. \ill{} 5$\frac{1}{2}$ stellæ polaris altitudo visa est partium 46.2.30. \ill{}.

Die 4 Jan. h. 5$\frac{1}{2}$ \ill{} \ill{} eiusd. stellæ altitudo supra polum in mer. par. 51.28$'$.45$''$.

Die 5 Jan. h. 5$\frac{1}{2}$ \ill{} matut. minima eiusd. stellæ altitudo p. 46.2.40. Sero eiusd. diei h. 5$\frac{1}{2}$ maxima altit. 51.28$'$.45$''$. Ac proinde \uncertain{consequens} altitudo veri poli fuit 48.45.37$\frac{1}{2}$. unde distantia stellæ a polo fuit 2.43$'$.7$\frac{1}{2}$$'''$ eadem quæ \uncertain{Tychoni} \ill{} \ill{} \ill{}.
\note{Le mot qui suit « h. » aux deux premières lignes, et celui qui suit « 5½ » à la deuxième, n'ont pas été lus ; le mot qui clôt la première ligne non plus. Les nombres se vérifient : la demi-somme de 51°28'45'' et de 46°2'30'' est 48°45'37½'', la demi-différence 2°43'7½''. Avec le 46.2.40 du 5 janvier, qui est écrit ainsi, elle ne tomberait pas juste. Le 7½ de la distance porte trois traits d'unité, comme écrit. La fin de la phrase, après « quæ », n'a pas été lue sauf un nom qui paraît être celui de Tycho.}

Vigesimo \uncertain{tertio} Decemb. anni 1620 Altitudo sol. merid. 17.45. quibus addita parall. 2$'$.58$''$ fiunt 17.47.58. His deinde addita declinationi 23.30.22. solis existentis in \add{Capricorne} 2.13. æquinoctialis altitudo erit 41.18.20.

25 Decembris eiusd. anni solis altitudo 17.48 $\odot$ in 4 gradu 16 \add{Capricorne} existenti, addita parall. fiunt 17.50.58. addita rursus declinat. 23.27.24. altitudo æquin. fit 41.18.22 maior propter refractiones \uncertain{non} observatas.
\note{Le signe rendu \add{Capricorne} a la forme d'un z à boucle ; le soleil est rendu $\odot$, comme écrit (un cercle). Les additions sont justes : 17°45' + 2'58'' = 17°47'58'', + 23°30'22'' = 41°18'20'' ; 17°48' + 2'58'' = 17°50'58'', + 23°27'24'' = 41°18'22''. Le mot devant « observatas » est bref et peu sûr.}

Au Reverend pere Denis Petau de la Société de Jesus

Au College de Clermont.
\note{adresse, de la même main, dans la moitié inférieure de la page, sur trois lignes puis une quatrième en retrait ; une tache de cire ou d'encre à droite de la dernière ligne.}

\note{Ce qui suit est la lettre montée sur l'onglet de la page de droite.}

Mon \uncertain{Reuerend} pere

Me souvenant de ce que vous me demandastes l'\uncertain{yver} \uncertain{passé} \uncertain{sur} la hauteur du pole pour latitude de Paris, il m'est souvenu de quelques observations qui \ill{} \ill{} au \ill{} de \uncertain{Jan.} 1621, qui sont \ill{} \ill{} \uncertain{suivant} lesquelles je vous envoye : si tant est que vous \ill{} \ill{} \ill{} : si vous ne les avez \uncertain{ou} je vous \uncertain{pourrois} \uncertain{servir} assez a propos. Je n'ay point \uncertain{retenu} de copie : c'est pourquoy je vous prie me renvoyer celle ci quand vous vous en serez servi. Je vous prie aussi de donner ce petit mot de lettre au R. pere \uncertain{Grandamy} \ill{} de ce qui est de la longitude de Paris s'il \uncertain{vous} \uncertain{vient} quelque chose de \ill{} \ill{} \ill{} fait \ill{} \uncertain{cependant} je me recommande a vos \uncertain{sainctes} prieres.
\note{Lettre d'une petite écriture rapide, très liée, dont beaucoup de mots ne se lisent qu'au sens ; ceux qui ne l'ont pas donné sont laissés \emph{illisibles}. Le millésime 1621 est net, et s'accorde avec les observations de janvier 1621 de la pièce précédente. Au-delà de « quelque chose de », la ligne porte un mot de six ou sept lettres, puis un nombre de quatre signes (1536 ? 1620 ?) que le passage ne permet pas de fixer. Le nom du jésuite à qui la lettre demande de transmettre un mot se lit « Grandamy » avec doute. Le cachet de la bibliothèque couvre la fin de la cinquième ligne.}

\uncertain{Vostre} tres \uncertain{affectionné} \ill{} Mar. Mers. M.
\note{souscription et signature, abrégées : « Mar. Mers. » suivi de « M. » pour Minime. Le premier mot est un grand signe bouclé qui peut abréger « Vostre », le suivant se lit « v. » ; « affect. » est net.}

Cette lettre est du P. Mersenne, minime
\note{d'une autre main, plus tardive, au bas du feuillet, à la suite de « R.D. 9555 ».}

\page{40}
\note{Même ouverture que la vue 39, le petit feuillet de la lettre étant rabattu vers la droite pour montrer son dos. La page de gauche est la même qu'à la vue 39. Le dos de la lettre, où transparaît l'écriture du recto, ne porte que l'adresse, écrite tête-bêche par rapport à la lettre :}

Au R. \ill{} pere Petau.
\note{adresse abrégée, d'une seule ligne ; entre « R. » et « pere », un mot bref et lié n'a pas été lu. Aucun cachet, aucune marque postale lisible.}

\end{document}
